\section{近期研究（2026 年 8--11 月）}
\paragraph{Gaussian-JEPA：预测高斯资产的隐空间结构}
2026.08~\cite{2608_15651}。同一对象可以被多组不同的高斯原语表示。Gaussian-JEPA 不要求从上下文精确重建被遮挡原语的原始参数，而用在线编码器读取可见令牌，预测由动量编码器生成的多个目标块表征。隐空间监督降低了对特定原语采样的一致性要求。

实验分别检验重采样、部分观测、冻结特征形状补全和部件分割等用途；重采样包含未参加预训练的 ModelNet40-GS 对象，并让 Gaussian-MAE 使用相同的重采样和下游协议。当前结论以每个对象固定 1000 个高斯为基础，不代表不同训练流程生成的两个资产天然等价，也未验证大规模动态场景。

\paragraph{LiteTex-GS：在二维高斯上按内容分配纹理容量}
2026.09~\cite{2609_23380}。带纹理的二维高斯可把几何和外观分开，但为所有基元固定高分辨率纹理会浪费容量。LiteTex-GS 用频率与误差线索逐级扩充局部纹理，按有效面积和渲染贡献管理基元，并按纹理分辨率调节更新幅度。这样几何原语数与每个原语的局部纹理细节可以分别分配。

作者在 Mip-NeRF 360 全部九个场景以及 Tanks and Temples 和 Deep Blending 做新视图对照。结果要与基准采用的图像分辨率与纹理设置一并解读。该论文解决的是带局部纹理的二维高斯训练，这组结果不代表可直接替换所有三维高斯渲染器。

\paragraph{SCION：让重复场景元素共享神经原语模板}
2026.10~\cite{2610_02322}。在含有重复结构的场景中，为每个实例完全独立保存高斯既占空间，也没留下可操作的对象位置。SCION 将场景拟合成可复用的原语词表和轻量的世界空间实例。训练分两层增减词表内的高斯和场景内的实例，由每个模板在不同实例中的观测共同约束；优化后可以在实例层移动或替换部件。

原文在 Mip-NeRF 360、Tanks and Temples、Deep Blending 和 Synthetic-NeRF 比较存储量相近时的图像质量，并用定性示例展示编辑。作者说明，真实重复弱或实例差异强时，模板容易欠拟合或退化成一次性表示；当前实例选择为固定的硬分配，两阶段优化和额外损失也增加训练成本。

\paragraph{直接保存条件切片所需的高维高斯参数}
2026.10~\cite{2610_22778}。传统高维高斯泼溅从联合协方差临时执行分块求逆与 Schur 补，才能在指定时间等查询变量下取三维切片。该方法为不透明度精度保存 Cholesky 因子，并以分解的耦合算子直接计算空间位置与尺度，减少渲染时的条件运算。

论文在五个静态及动态基准和高斯、Beta 两类核中作成对比较，并分别报告切片与整帧加速。局部切片快不能推出光栅化或端到端同倍加速，且配对方法共享原语预算时实际基元数仍可不同。更大场景的内存、可见性和切块瓶颈，以及其他核函数的适用性仍待验证。

\paragraph{按实际合成贡献而非投影次数累计增密信号}
2026.10~\cite{2610_27108}。传统高斯增密统计可能把被前景遮住的高斯也计入视图数，令少量真正可见的梯度被稀释。该方法先使完全遮挡的原语不再污染统计，再用各视图中实际合成贡献的累积平均量归一化屏幕空间梯度，兼顾不同深度和投影范围。

主文在 StopThePop 上检查重遮挡数据的原语数与感知画质，补充材料还将两项改动移植到 2DGS 和原始 3DGS，并对 Mip-NeRF 360 和不同训练长度报告结果。遮挡统计依赖渲染器保留对合成贡献的计算，其他近似合成路径的收益仍需按实现对照。

\paragraph{用多视图重建损失的一阶变化决定高斯增删}
2026.10~\cite{2610_34344}。固定预算训练时，原语分配通常用不透明度或位置梯度作为代理。该研究在已有原语位置引入辅助质量变量，计算多视图损失对其局部变化的敏感度，在统一预算下给高效候选增加容量，并减少低效原语。实际产生新原语的方式仍交给既有训练后端。

主文固定渲染器、优化器和 MCMC 转移，在近似相同原语数量下比较控制指标。Mip-NeRF 360 主结果使用七个场景，排除了 Flowers 和 Treehill。原文也将评分移植到常规克隆、分裂和删除，但已有位置的局部导数只能指示容量送到哪里，不能保证全局结构最优。

\paragraph{Gaussian Splatoms：用各向同性高斯直接表达电子断层中的原子}
2026.11~\cite{2611_23444}。Gaussian Splatoms 给每个原子分配一个各向同性三维高斯，通过投影拟合电子断层观测并学习位置、振幅和尺度。原子间距和邻接尺度先验限制过近或不合理的解，投影几何也可调整。这一模型针对电子束投影，与普通 RGB 高斯的遮挡合成用途不同。

实测 FePt 数据包含 68 幅约正负 65 度的投影。论文对颗粒内部原子与人工筛订的 GENFIRE 参考比较，报告真阳性率 98.3\%、假阳性率 0.42\%、位置均方根差 0.32 埃。参考来自既有重建而非独立测量真值，这组内部原子指标也不能代表边界原子的准确率。
